\section{The tail of the kernel}\label{s:tail}

In this appendix we prove the bound on $\rho(M_1)$ in Theorem~\ref{t:inputs}: at $(\cs, \ct) = (0.961, 0.935)$, with $\bar\alpha = 0.18094$ and $\bar\beta = 0.36751$, the kernel \eqref{e:kernel} satisfies $\rho(M_1) \le 0.98779$. The trails of length at most $20$ are enumerated exactly. The longer trails are followed exactly for $20$ steps and are then dominated by a weighted walk with a memory of $10$ steps on a periodic quotient of $\Z^2$, and the walk after $260$ steps is bounded by a sub-invariant vector. The point of the quotient is that the walk is still required to end at the site $v$, or at one of its images, so that the bound keeps the closure of the trails. A bound by open trails would be useless here, since the open-trail growth rate $q \approx 0.98$ of (H2) is close to $1$.

\subsection{Reduction to two scalar sums}\label{s:tailred}

Throughout, $x_{\mathrm s} = 1/\cs$, $x_{\mathrm t} = 1/\ct$, and $h = (h_B, h_C) = (1, h_C)$, where $h_C$ is the double-precision number nearest to $1.23$ (the programs below use this number, and any positive vector is admissible). For a first-visit passage $i$ of a trail let $\kappa_h(i) = \sum_{s \ge 0} \bigl( \kappa_{B, s}(i) + h_C \kappa_{C, s}(i) \bigr)$, with $\kappa_{Y, s}$ as in \S\ref{s:kernel}. For $X \in \{B, C\}$ and a set $\mathcal T \subset \mathcal T^X_n$ let
\begin{equation}\label{e:mT}
\mathsf m(\mathcal T) = \sum_{w \in \mathcal T} 3^{-n} \Bigl( \frac35 \Bigr)^{\mathrm{rev}(w)} \sum_k \kappa_h(\sigma_w(k)) \prod_{l \le k} x_{\sigma_w(l)}, \qquad \mathsf m_X(n) = \mathsf m(\mathcal T^X_n).
\end{equation}
By \eqref{e:kernel}, $(M_1 h)_X = \sum_{n \ge 1} \mathsf m_X(n)$. Since $|u - v|_1 = 2$, a trail from $u$ to $v$ has even length, so $\mathsf m_X(n) = 0$ for odd $n$.

\begin{lemma}\label{l:CW}
Let $M$ be a $2 \times 2$ matrix with nonnegative entries, let $h \in (0, \infty)^2$, and let $\mu = \max_X (M h)_X / h_X$. Then $\rho(M) \le \mu$.
\end{lemma}

\begin{proof}
We have $M h \le \mu h$ entrywise, hence $M^k h \le \mu^k h$, so every entry of $M^k$ is at most $\mu^k \max_Y h_Y / \min_Y h_Y$, and $\rho(M) = \lim_k \|M^k\|^{1/k} \le \mu$.
\end{proof}

\emph{The orderings.} For a trail of length $n \le 20$, $\sigma_w$ lists the first-visit passages $i$ in increasing order of $(x_i - 1)/(x_i \kappa_h(i))$ (the value $+\infty$ if $\kappa_h(i) = 0$). Exchanging two consecutive passages $i, j$ with prefix product $\Pi$ changes their contribution from $\Pi (x_i \kappa_h(i) + x_i x_j \kappa_h(j))$ to $\Pi (x_j \kappa_h(j) + x_j x_i \kappa_h(i))$; hence this order minimizes the $w$-term of \eqref{e:mT}, and ties do not affect its value. For a trail of length $n \ge 22$, $\sigma_w$ lists the straight first visits in path order and then the turning first visits in path order (\emph{straights first}). Both choices are functions of $w$ alone and do not depend on $\theta$, as Lemma~\ref{l:Mtheta} requires.

\emph{Kinds and phases.} The \emph{kind} of a passage is $\mathrm s$ (straight first visit), $\mathrm t$ (turning first visit) or $\mathrm r$ (revisit). With the straights-first order, the $w$-term of \eqref{e:mT} is a sum over the responding first visit $j$ of a product of one factor per passage. If $j$ is a turn, the passages before $j$ have the factors $x_{\mathrm s}/3$, $x_{\mathrm t}/3$, $1/5$ according to their kind, the passage $j$ has the factor $x_{\mathrm t} \kappa_h(j)/3$, and the passages after $j$ have the factors $x_{\mathrm s}/3$, $1/3$, $1/5$; indeed the first visits that precede $j$ in $\sigma_w$ are all straight first visits and the turning ones up to $j$. If $j$ is straight, the factors are $x_{\mathrm s}/3$, $1/3$, $1/5$ before $j$, the factor $x_{\mathrm s} \kappa_h(j)/3$ at $j$, and $1/3$, $1/3$, $1/5$ after $j$. We record this with four \emph{phases}: $\mathrm A$ (before a turning response), $\mathrm C$ (after it), $\mathrm B$ (before a straight response), $\mathrm D$ (after it). A passage of kind $\mathrm k$ in phase $\varphi$ has the weight $\varpi_\varphi(\mathrm k)$ given by
\begin{equation}\label{e:phaseweights}
\begin{array}{c|ccc}
 & \mathrm s & \mathrm t & \mathrm r \\ \hline
\mathrm A & x_{\mathrm s}/3 & x_{\mathrm t}/3 & 1/5 \\
\mathrm B & x_{\mathrm s}/3 & 1/3 & 1/5 \\
\mathrm C & x_{\mathrm s}/3 & 1/3 & 1/5 \\
\mathrm D & 1/3 & 1/3 & 1/5
\end{array}
\end{equation}
and the responding passage moves the phase from $\mathrm A$ to $\mathrm C$ (kind $\mathrm t$, weight $x_{\mathrm t} \kappa_h/3$) or from $\mathrm B$ to $\mathrm D$ (kind $\mathrm s$, weight $x_{\mathrm s} \kappa_h/3$). The $w$-term of \eqref{e:mT} is the sum, over the phase sequences that start in $\mathrm A$ or $\mathrm B$ and switch exactly once, of the products of these weights.

\subsection{The memory walk on a quotient}\label{s:memwalk}

Let $m = 10$ and $\Lambda = \{ (28(a + b), 28(a - b)) : a, b \in \Z \}$, and let $\mathsf Q = \Z^2/\Lambda$, a set of $1568$ classes; $[z]$ denotes the class of $z \in \Z^2$. For $\lambda = (28(a + b), 28(a - b)) \neq 0$ we have $|\lambda|_1 = 56 \max(|a|, |b|) \ge P := 56$.

A \emph{window} is a sequence $\omega = (d_1, \dots, d_m)$ of unit vectors such that the path $z_0 = 0$, $z_k = z_{k - 1} + d_k$ is a trail; let $\Omega$ be the set of windows ($|\Omega| = 60{,}092$). A unit vector $d$ is an \emph{admissible step} from $\omega$ if $d \neq -d_m$ and $\{z_m, z_m + d\}$ is not a bond of the path; then $\omega \cdot d = (d_2, \dots, d_m, d) \in \Omega$, and the kind of the step is $\mathrm r$ if $z_m \in \{z_0, \dots, z_{m - 1}\}$, $\mathrm s$ if $d = d_m$, and $\mathrm t$ otherwise. Let $\kappa^\circ(\omega, d)$ be the value of $\kappa_h$ for a first-visit passage with arrival direction $d_m$ and out-direction $d$ when all three pairs are live: $\bar\alpha + 2 \bar\beta h_C$ for a straight step and $\bar\beta + (\bar\alpha + \bar\beta) h_C$ for a turn.

We write $\mathsf E = (1, 0)$, $\mathsf N = (0, 1)$, $\mathsf W = (-1, 0)$, $\mathsf S = (0, -1)$ for the four directions. Fix $X \in \{B, C\}$, a first direction $D_0 \in \{\mathsf N, \mathsf W, \mathsf S\}$ or $D_0 = \ast$ (no restriction), and a time $t_* \ge 20$. The \emph{memory walk} is the sequence of functions $V_t \colon \Omega \times \mathsf Q \times \{\mathrm A, \mathrm B, \mathrm C, \mathrm D\} \to [0, \infty)$, $t \ge 20$, and of numbers $r(n)$, $n > 20$, defined as follows.
\begin{enumerate}
\item[(I)] $V_{20}(\omega, \xi, \varphi)$ is the sum, over the trails $w$ of $\Z^2 \setminus \{c\}$ of length $20$ from $u$, with first step $D_0$ (if $D_0 \neq \ast$), that do not visit $v$, whose last $m$ steps form $\omega$ and whose end lies in $\xi$, of the exact weight of $w$ in phase $\varphi$: for $\varphi \in \{\mathrm A, \mathrm B\}$ the product of the weights $\varpi_\varphi$ of its $20$ passages, and for $\varphi = \mathrm C$ (respectively $\mathrm D$) the sum over its turning (straight) first visits $j$ of the products with the switch at $j$, all with the exact kinds and the exact values $\kappa_h(j)$ of \S\ref{s:kernel}.
\item[(II)] For $t \ge 20$, $V_{t + 1}(\omega \cdot d, \xi + d, \varphi')$ is the sum over $\omega$, $\varphi$ and the admissible steps $d$ of $V_t(\omega, \xi, \varphi)$ times the weight of the step: $\varpi_\varphi(\mathrm k)$ if $\varphi' = \varphi$, the response weight $x_{\mathrm k} \kappa^\circ(\omega, d)/3$ if $(\varphi, \varphi') = (\mathrm A, \mathrm C)$ and $\mathrm k = \mathrm t$ or $(\varphi, \varphi') = (\mathrm B, \mathrm D)$ and $\mathrm k = \mathrm s$, and $0$ otherwise; here $\mathrm k$ is the kind of the step. If $t < t_*$, the following rules modify the step.
\begin{enumerate}
\item[(R1)] If $\xi \in \{[u], [u + D_0]\}$ (only $\xi = [u]$ if $D_0 = \ast$): a step along the bond $\{u, u + D_0\}$, that is, from $[u]$ in direction $D_0$ or from $[u + D_0]$ in direction $-D_0$, has weight $0$; a step of kind $\mathrm r$ has weight $0$ if $D_0 \neq \ast$ and weight $1/5$ (without phase change) if $D_0 = \ast$; every other step has weight $1/5$ and no phase change.
\item[(R2)] In the response weight, $\kappa^\circ(\omega, d)$ is replaced by the sum of the coefficients, weighted by $h$, of the live pairs, where a pair is now dead if it contains an element of $[c]$.
\end{enumerate}
\item[(III)] $r(t + 1) = \sum_\omega \bigl( V_{t + 1}(\omega, [v], \mathrm C) + V_{t + 1}(\omega, [v], \mathrm D) \bigr)$. If $t + 1 \le t_*$, after this is recorded $V_{t + 1}(\cdot, [v], \cdot)$ and $V_{t + 1}(\cdot, [c], \cdot)$ are set to $0$.
\end{enumerate}
\emph{Mode A} is $t_* = 109$ with $D_0 \in \{\mathsf N, \mathsf W, \mathsf S\}$; \emph{mode B} is $t_* = 53$ with $D_0 = \ast$. We write $r^{\mathrm A}_{X, D_0}(n)$ and $r^{\mathrm B}_X(n)$. Let $\mathcal T^{X, D_0}_n$ be the set of trails of $\mathcal T^X_n$ with first step $D_0$ (all of $\mathcal T^X_n$ if $D_0 = \ast$); the first step of a trail of $\mathcal T^X_n$ is not $\mathsf E$, since $u + \mathsf E = c$.

\begin{lemma}\label{l:memdom}
Let $X \in \{B, C\}$, with the straights-first orderings.
\begin{enumerate}
\item In mode A, $\mathsf m(\mathcal T^{X, D_0}_n) \le r^{\mathrm A}_{X, D_0}(n)$ for $20 < n \le 2P - 5 = 107$.
\item In mode B, $\mathsf m_X(n) \le r^{\mathrm B}_X(n)$ for every $n > 53$.
\end{enumerate}
\end{lemma}

\begin{proof}
We first collect the geometric facts. Let $w = (v_0, \dots, v_n) \in \mathcal T^{X, D_0}_n$ and $0 \le t \le n$. Then $t \ge |v_t - u|_1$ and $n - t \ge |v_t - v|_1$. For $\lambda \in \Lambda \setminus \{0\}$, $|c + \lambda - u|_1 \ge P - 1$ and $|c + \lambda - v|_1 \ge P - 1$; $|v + \lambda - u|_1 \ge P - 2$ and $|v + \lambda - v|_1 \ge P$; $|u + \lambda - u|_1 \ge P$ and $|u + \lambda - v|_1 \ge P - 2$; $|u + D_0 + \lambda - u|_1 \ge P - 1$ and $|u + D_0 + \lambda - v|_1 \ge P - 3$, since $|u + D_0 - v|_1 \le 3$. Consequently, if $v_t \in [c] \cup [v] \cup [u] \cup [u + D_0]$, or $v_t$ is a neighbour of an element of $[c]$, but $v_t$ is not $c$, $v$, $u$, $u + D_0$ or a neighbour of $c$ respectively, then $n \ge 2P - 4$, and moreover $t \ge P - 2$. In mode A ($n \le 2P - 5$), therefore, every element of $[c]$, $[v]$, $[u]$, $[u + D_0]$ that $w$ meets, or meets a neighbour of, is $c$, $v$, $u$, $u + D_0$ itself. In mode B the rules act only at times $t \le 53 = P - 3$, when the same conclusion holds for every $n$. Finally, the admissibility and the kind of a step depend only on the window, and they are computed in $\Z^2$, without reference to $\Lambda$.

Now fix $w$ and a responding first visit $j$, and let $\pi(w, j)$ be the sequence of states $(\omega_t, [v_t], \varphi_t)$, $20 \le t \le n$, where $\omega_t$ is the sequence of the last $m$ steps of $w$ before time $t$ and $\varphi_t$ is the phase after $t$ passages. We compare the factor of passage $t$, $20 \le t < n$, in the $(w, j)$-term of the phase representation with the weight of the step of $\pi(w, j)$ from time $t$ to $t + 1$.
\begin{itemize}
\item The step is admissible: $w$ is a trail, so it never reverses and never reuses a bond of the window. The chain sets the mass at $[c]$ to $0$ only when $w$ is at $c$, which never happens, and records and deletes the mass at $[v]$ only at a time $t \le n$ when $v_t = v$, which happens only at $t = n$, after $r(n)$ is recorded.
\item If the kind of the step is $\mathrm r$, then $v_t$ was visited by $w$ at a time in $[t - m, t - 1]$, so passage $t$ is a revisit of $w$, with the same weight $1/5$ and no response.
\item If the kind is $\mathrm s$ or $\mathrm t$ and passage $t$ is a first visit of $w$, the kind is its true type and the weights agree. If the kind is $\mathrm s$ or $\mathrm t$ and passage $t$ is a revisit of $w$ (of a site visited before time $t - m$), the true weight is $1/5$ and the chain weight is at least $1/3$; such a passage is not the responding one.
\item At the responding passage $t = j \ge 20$, which is a first visit, the chain weight is $x \kappa^\circ/3$, or $x$ times the value of (R2) over $3$; the pairs dropped by (R2) contain an element of $[c]$ adjacent to $v_j$, which is $c$ itself by the geometric facts, so they are dead in the kernel. Hence the chain weight is at least $x \kappa_h(j)/3$.
\item Rule (R1) acts at $v_t \in [u] \cup [u + D_0]$, that is (geometric facts), at $v_t = u$ or $v_t = u + D_0$, which $w$ visited at times $0$ and $1$. Passage $t \ge 20$ is then a revisit of $w$, with weight $1/5$ and no response, and $w$ does not use the bond $\{u, u + D_0\}$ again. A step of kind $\mathrm r$ at $u$ or $u + D_0$ would mean a third visit at a time in $[t - m, t - 1] \subset [10, t - 1]$, and $w$ cannot make it: $u$ has three bonds in $\Z^2 \setminus \{c\}$, one used at time $0$, and $u + D_0 \neq v$ has four, two used at time $1$.
\end{itemize}
Hence the weight of $\pi(w, j)$ from time $20$ to time $n$ is at least the product of the factors of the passages $20, \dots, n - 1$ in the $(w, j)$-term, the path is never deleted before time $n$, and at time $n$ it is at $([v], \varphi_n)$ with $\varphi_n \in \{\mathrm C, \mathrm D\}$. The factors of the passages $0, \dots, 19$ form the exact weight of the prefix $(v_0, \dots, v_{20})$ in phase $\varphi_{20}$, with the switch at $j$ if $j < 20$. For fixed states $\pi$ the pairs $(w, j)$ with $\pi(w, j) = \pi$ differ only in their prefix of length $20$ and, if $j < 20$, in $j$, and the corresponding prefix weights are among the terms of $V_{20}$ at the initial state of $\pi$, by (I). Since all weights are nonnegative, summing over $(w, j)$ gives $\mathsf m(\mathcal T^{X, D_0}_n) \le r(n)$.
\end{proof}

\begin{lemma}\label{l:remainder}
Let $V_t$ be a memory walk with $t_* < T$, and let $\mathbb W$ be its weight matrix on $\Omega \times \{\mathrm A, \mathrm B, \mathrm C, \mathrm D\}$ without rules, $\mathbb W((\omega, \varphi), (\omega \cdot d, \varphi'))$ the weight of the step $d$ from phase $\varphi$ to phase $\varphi'$ in (II). Suppose that $z \colon \Omega \times \{\mathrm A, \mathrm B, \mathrm C, \mathrm D\} \to [0, \infty)$, $\lambda \in (0, 1)$ and $a \ge 0$ satisfy $z \mathbb W \le \lambda z$ and $V_T(\omega, \xi, \varphi) \le a \, z(\omega, \varphi)$ for all $\omega, \xi, \varphi$. Then
\begin{equation*}
\sum_{n > T} r(n) \le \frac{a \lambda}{1 - \lambda} \sum_{\omega \in \Omega} \bigl( z(\omega, \mathrm C) + z(\omega, \mathrm D) \bigr).
\end{equation*}
\end{lemma}

\begin{proof}
For $t \ge T > t_*$ no rule acts, and the predecessor of a state $(\omega', \xi', \varphi')$ has the site $\xi' - d'_m$, where $d'_m$ is the last step of $\omega'$. Hence, if $V_t \le a \lambda^{t - T} z$ on all states, then
\begin{equation*}
V_{t + 1}(\omega', \xi', \varphi') = \sum_{\omega, \varphi} V_t(\omega, \xi' - d'_m, \varphi) \, \mathbb W((\omega, \varphi), (\omega', \varphi')) \le a \lambda^{t - T} (z \mathbb W)(\omega', \varphi') \le a \lambda^{t + 1 - T} z(\omega', \varphi'),
\end{equation*}
where the sum runs over the $\omega$ with $\omega \cdot d'_m = \omega'$. By induction this holds for all $t \ge T$, and summing $r(t) \le a \lambda^{t - T} \sum_\omega (z(\omega, \mathrm C) + z(\omega, \mathrm D))$ over $t > T$ gives the claim.
\end{proof}

The matrix $\mathbb W$ is block triangular in the phases ($\mathrm A \to \mathrm A, \mathrm C$; $\mathrm B \to \mathrm B, \mathrm D$; $\mathrm C \to \mathrm C$; $\mathrm D \to \mathrm D$), and its spectral radius is that of the block $\mathrm A \to \mathrm A$, the memory-$10$ walk with the weights $x_{\mathrm s}/3$, $x_{\mathrm t}/3$, $1/5$. A nonnegative candidate $z$ is obtained by power iteration in the diagonal blocks and approximate solution of
$z_{\mathrm C}(\lambda-\mathbb W_{\mathrm{CC}})=z_{\mathrm A}\mathbb W_{\mathrm{AC}}$ and
$z_{\mathrm D}(\lambda-\mathbb W_{\mathrm{DD}})=z_{\mathrm B}\mathbb W_{\mathrm{BD}}$;
only the inequality $z\mathbb W\le\lambda z$ is used, and it is checked directly.
The constructed vector is zero in each phase at exactly the $64$ windows of indegree zero, and positive at every other window.
No update can put mass at a window of indegree zero, so $V_T$ vanishes there for $T\ge21$.
In applying Lemma~\ref{l:remainder}, we check this zero-support condition and take ratios only where $z>0$.

\subsection{The certified numbers}\label{s:tailnumbers}

\begin{proposition}\label{p:tail}
At $(\cs, \ct) = (0.961, 0.935)$, $\bar\alpha = 0.18094$, $\bar\beta = 0.36751$ and with the orderings of \S\ref{s:tailred}, the entries of Table~\ref{tab:tail} are upper bounds for the sums in its first column, and $(M_1 h)_B / h_B \le 0.987763$, $(M_1 h)_C / h_C \le 0.968624$. Hence $\rho(M_1) \le 0.98777$, which implies the bound $0.98779$ of Theorem~\ref{t:inputs}.
\end{proposition}

\begin{table}[ht]
\centering
\begin{tabular}{lcc}
 & row $B$ & row $C$ \\ \hline
$\sum_{n \le 20} \mathsf m_X(n)$ (exact enumeration) & $0.594580$ & $0.841681$ \\
$\mathsf m_X(22) \le \sum_{D_0} r^{\mathrm A}_{X, D_0}(22)$ & $0.026669$ & $0.022644$ \\
$\sum_{24 \le n \le 106} \mathsf m_X(n) \le \sum_{D_0} \sum_n r^{\mathrm A}_{X, D_0}(n)$ & $0.302932$ & $0.268029$ \\
$\sum_{108 \le n \le 260} \mathsf m_X(n) \le \sum_n r^{\mathrm B}_X(n)$ & $0.057384$ & $0.053319$ \\
$\sum_{n > 260} \mathsf m_X(n)$ (Lemma~\ref{l:remainder}) & $0.006199$ & $0.005737$ \\ \hline
$(M_1 h)_X$ & $0.987763$ & $1.191407$ \\
$(M_1 h)_X / h_X$ & $0.987763$ & $0.968624$
\end{tabular}
\caption{The certified bound at $(\cs, \ct) = (0.961, 0.935)$, $\bar\alpha = 0.18094$, $\bar\beta = 0.36751$, $h = (1, 1.23)$. Every entry is the exact sum of the upward-rounded outputs, rounded up in the last digit; the two last rows are rounded up from the exact sums of these outputs, not from the rounded entries. In row $B$ the class $D_0 = \mathsf S$ is the mirror image of the class $\mathsf N$ under $y \mapsto -y$, which fixes $c$, $u$, $v$ and $\Lambda$ and preserves kinds, pair types and dead pairs, so it is counted as the class $\mathsf N$.}\label{tab:tail}
\end{table}

\begin{proof}
By Lemma~\ref{l:CW} it suffices to bound $(M_1 h)_X = \sum_n \mathsf m_X(n)$ for both rows, and we split the sum over the even $n$ as in Table~\ref{tab:tail}. The trails of length at most $20$ are enumerated exactly. For $22 \le n \le 106$ we use Lemma~\ref{l:memdom}(1) and $\mathcal T^X_n = \bigcup_{D_0} \mathcal T^{X, D_0}_n$, a disjoint union. For $108 \le n \le 260$ we use Lemma~\ref{l:memdom}(2); note that $n = 108$ is beyond the range of mode A, because the rules (R1) at $[u + D_0]$ and (R2) can act at images for closed trails of length $2P - 4 = 108$. For $n>260$ we use Lemma~\ref{l:remainder} with $T=260$ for the walk of mode B, taking $\lambda$ to be the stored binary64 value \texttt{lcheck} used in both the sub-invariance check and the remainder calculation, with $0<\lambda\le0.979532789<0.97953279$.
The nonnegative vector $z$ satisfies the directly checked inequality $z\mathbb W\le\lambda z$.
We verify $V_{260}(\omega,\xi,\varphi)=0$ whenever $z(\omega,\varphi)=0$, and choose $a$ as an upward-rounded upper bound for
$\max_{z(\omega,\varphi)>0}V_{260}(\omega,\xi,\varphi)/z(\omega,\varphi)$. The entries of the table are computed as described in \S\ref{s:tailtrust}.
\end{proof}

\subsection{What is computed, and what must be trusted}\label{s:tailtrust}

\emph{Programs.} \texttt{tailmem5.c} computes, for one row and one class $D_0$, the exact sums $\mathsf m(\mathcal T^{X, D_0}_n)$, $n \le 20$, by depth-first search with the per-trail optimal ordering, the initial vector (I), and the memory walk of mode A up to $t = 109$, printing $r^{\mathrm A}_{X, D_0}(n)$ for every $n$. \texttt{tailmem4.c} does the same without the class split for mode B up to $T = 260$, computes $z$ by power iteration, checks $z \mathbb W \le \lambda z$, and prints $a$, $\lambda$ and the bound of Lemma~\ref{l:remainder}.
The archived mode-A summary through $109$ is a memory-walk mass, and includes the unproved-for-trails value at $108$; only its per-length outputs through $106$ are used in Table~\ref{tab:tail}.
The unsplit mode-B production runs use \texttt{tailmem4.c}.
The archived unused unsplit option of \texttt{tailmem5.c} additionally deletes a window-detected revisit at $[u]$ while the rules act, whereas (R1) with $D_0=\ast$ retains it with weight $1/5$.
This additional deletion cannot remove a genuine trail in the valid range: the known time-$0$ visit, a window-detected later visit, and another departure would exceed the three available bonds at $u$.
It defines a stronger auxiliary rule and is not used for the mode-B entries of the table.
The obsolete validity comments and summary labels in the archived sources must not be used in place of Lemma~\ref{l:memdom}. The table is assembled from the five files
\texttt{r961A\_B\_1.txt}, \texttt{r961A\_B\_2.txt},
\texttt{r961A\_C\_1.txt}, \texttt{r961A\_C\_2.txt},
\texttt{r961A\_C\_3.txt} and the two files
\texttt{r961B\_B.txt}, \texttt{r961B\_C.txt}, archived in
\texttt{agents/75-audit-tail/out/rerun/} in the full project.
The commands are recorded in \texttt{agents/75-audit-tail/code/run\_all.sh}
and \texttt{run\_rest1.sh}, \texttt{run\_rest2.sh}; the rerun used
\texttt{gcc -O2 -frounding-math -ffp-contract=off}.
The exact assembly is implemented in \texttt{agents/75-audit-tail/code/assemble75.py}:
sum the printed decimals as rational numbers, using the class multiplicities in the caption,
the mode-A values for $22\le n\le106$, the mode-B values for $108\le n\le260$,
and the printed mode-B remainder, without subtracting any rounded tail totals.
The denominator is the exact binary64 number $h_C=2769713770832855/2251799813685248$.
The assembled numerators are $0.987762443170525$ and $1.1914067844431$;
their quotients are at most $0.987763$ and $0.968624$. The state space of the walk has $60{,}092 \times 1{,}568 \times 4 \approx 3.8 \cdot 10^8$ entries (three gigabytes in single precision); each of the five mode-A runs takes about five minutes and each of the two mode-B runs about twenty minutes on one core.

\emph{Arithmetic.}
The enumeration and walk bounds are evaluated as sums of products of nonnegative quantities.
The constants $\cs,\ct$ are decreased and $\bar\alpha,\bar\beta$ increased by one unit in the last place after decimal conversion;
$1/3$, $1/5$, $1/\cs$ and $1/\ct$ are computed by upward division.
The walk arrays are stored in binary32 and the enumeration in binary64.
Bound-producing products, sums and conversions, and their output decimals, are rounded upward, with gradual underflow.
The ordering search and the construction of a candidate $z$ need not use upward rounding: any chosen ordering gives an upper bound for the per-trail minimum, and the stored candidate is subsequently treated as an exact dyadic vector.
In checking $z\mathbb W\le\lambda z$, the left side is rounded upward and the right side downward; $1-\lambda$ is rounded downward in the remainder formula.
Acceptance requires finite nonnegative data, $0<\lambda<1$, the checked sub-invariance inequality, and $V_T=0$ wherever $z=0$.
Ratios defining $a$ are taken only at positive coordinates of $z$, with upward rounding; nonfinite values or failed checks reject the certificate.
The table uses exact rational sums of the upward-printed per-length bounds and the upward-printed remainder, followed by decimal ceilings.
The programs must be compiled with \texttt{-frounding-math}, without unsafe reassociation; the archived rerun also used \texttt{-ffp-contract=off}.
Correct directed arithmetic, conversions and decimal output, and gradual underflow, are part of the computational trust assumption.
The explicit support and finiteness guards are implemented in \texttt{tailmem4\_checked.c} and \texttt{tailmem5\_checked.c}, archived under \texttt{reports/ADJ-R5/audit/} relative to the project root; the production records identified above were generated by the original sources. Their zero-support condition follows from the indegree argument above.

\emph{Independent checks.} (1) An independent enumeration from the definition \eqref{e:kernel} reproduces, for every $n \le 20$ and every class, the exact sums used in the table to a relative error of $7 \cdot 10^{-13}$; it also reproduces, for every $n \le 22$, the kernel with the forward ordering computed by a third, symbolic enumeration (to $4 \cdot 10^{-10}$, the precision of its output), and for $n \le 16$ the trail counts and masses of a fourth implementation. (2) An independent implementation of the memory walk of \S\ref{s:memwalk}, run at memory $m = 4$ and $m = 6$ on the quotient by $\langle (9, 9), (9, -9) \rangle$, reproduces the output of \texttt{tailmem5.c} at the same parameters to a relative error of $10^{-5}$ for every $n$, both rows, all classes and both modes. (3) The exact sums $\mathsf m(\mathcal T^{X, D_0}_n)$ with the straights-first ordering were enumerated up to $n = 26$; they are dominated by the memory walk at every $n$ and in every class, for $m = 4, 6$ and for $m = 10$. At $m = 10$ the ratio of the bound to the exact sum lies, in the five classes, between $1.03$ and $1.21$ at $n = 22$ and between $1.19$ and $1.81$ at $n = 26$; for the rows it is $1.05$ ($B$) and $1.08$ ($C$) at $n = 22$. (4) The numbers of Table~\ref{tab:tail} come from the archived recomputation with
$\bar\alpha=0.18094$, $\bar\beta=0.36751$ entered directly, as identified above.
Exact assembly of the earlier outputs with $0.18093$, $0.36750$ gives the row bounds
$0.987731$ and $0.968592$.
Every printed per-$n$ output of the two runs differs by a factor in
$[1.0000322,1.0000333]$.
Since every term of \eqref{e:mT} carries exactly one response coefficient, linearly and with a nonnegative weight,
and the earlier orderings may be kept, the earlier run alone also gives
$\rho(M_1)\le0.987731\cdot0.18094/0.18093<0.98779$.

\emph{What must be trusted.} A reader who accepts the lemmas of this appendix has to trust that the implementations and explicit checks identified above implement the walk of \S\ref{s:memwalk}, the exact enumeration and the support-aware check of Lemma~\ref{l:remainder}, and that the floating-point environment implements IEEE~754 arithmetic with directed rounding. The checks (1)--(3) test the first point independently; they are not themselves rigorous.

\begin{remark}\label{r:tailtypedep}
With the type-dependent coefficients of Remark~\ref{r:typedep} in place of $\bar\alpha, \bar\beta$ (and the coefficient $\beta'$ of the surviving pair when exactly one of $f, g$ is known to be removed), a variant of the same programs (\texttt{tailmem6.c}) with $h_C = 1.20$ gives $(M_1 h)_B \le 0.91301$ and $(M_1 h)_C / h_C \le 0.91825$, hence $\rho(M_1) \le 0.91825$. The six coefficient bounds of Remark~\ref{r:typedep} were confirmed by an independent interval branch and bound, and the exact part $n \le 20$ of this computation by the independent enumeration; the memory-walk part was not recomputed. The seven production records for this variant are
\texttt{c961\_935\_A\_B\_1.txt}, \texttt{c961\_935\_A\_B\_2.txt},
\texttt{c961\_935\_A\_C\_1.txt}, \texttt{c961\_935\_A\_C\_2.txt},
\texttt{c961\_935\_A\_C\_3.txt}, \texttt{c961\_935\_B\_B.txt},
\texttt{c961\_\allowbreak 935\_\allowbreak B\_\allowbreak C.txt} in \texttt{agents/\allowbreak 62-kernel-tail/\allowbreak out/}.
Their exact assembly, using the binary64 value of $1.20$, gives row bounds
$0.913009554182608$ and a quotient less than $0.918245584$.
The unsplit implementation uses the additional start-site deletion described in \S\ref{s:tailtrust}, justified by the same degree argument, and the remainder uses the nonnegative-vector formulation of Lemma~\ref{l:remainder}.
This variant is not used in the main line.
\end{remark}
